\clearpage
\section{Hopfieldův model}\label{sec:ai-hopfield}

Dopředná síť neobsahuje orientovaný cyklus: informace postupuje od vstupu k~výstupu a~výpočet skončí. \emph{Rekurentní} síť naopak dovoluje zpětné spoje. Její nový stav proto závisí na stavu předchozím a~síť může uchovávat informaci v~čase.

\begin{figure}[ht]
    \centering
    \begin{tikzpicture}[
    neuron/.style={circle, draw=darkblue, fill=darkblue!10, minimum size=9mm},
    conn/.style={-Latex, thick, draw=darkblue}
]
    \node[neuron] (f1) at (0,0) {};
    \node[neuron] (f2) at (1.6,0) {};
    \node[neuron] (f3) at (3.2,0) {};
    \draw[conn] (f1) -- (f2);
    \draw[conn] (f2) -- (f3);
    \node[above=7mm of f2] {dopředná síť};

    \node[neuron] (r1) at (5.5,0) {};
    \node[neuron] (r2) at (7.1,0) {};
    \node[neuron] (r3) at (8.7,0) {};
    \draw[conn] (r1) -- (r2);
    \draw[conn] (r2) -- (r3);
    \draw[conn] (r3) to[bend left=55] node[above] {zpětný spoj} (r1);
    \node[above=13mm of r2] {rekurentní síť};
\end{tikzpicture}


    \caption{Rozdíl mezi dopřednou a~rekurentní sítí.}
    \label{fig:ai-dopredna-rekurentni}
\end{figure}

\subsection{Definice a~dynamika}

\begin{definition}\label{def:ai-hopfield}
    \emph{Hopfieldova síť} s~\(n\) neurony je dvojice
    \[
        \mathcal H_n=(N,w),
    \]
    kde \(N=\set{1,\ldots,n}\) je množina neuronů a
    \[
        w\colon N\times N\to\R
    \]
    je symetrická váhová funkce splňující
    \[
        w(i,j)=w(j,i),
        \qquad
        w(i,i)=0.
    \]
\end{definition}

Váhu \(w(i,j)\) budeme krátce zapisovat \(w_{ij}\). Každému neuronu \(i\) navíc přiřadíme práh \(\theta_i\in\R\). Stav sítě v~čase \(t\) je vektor
\[
    \mathbf{s}^{(t)}
    =
    \left(
        s_1^{(t)},\ldots,s_n^{(t)}
    \right)
    \in\set{-1,1}^{n}.
\]
Síť je plně propojená: každý neuron přijímá stav všech ostatních neuronů, vynásobí je vahami a~vytvoří lokální pole
\[
    h_i^{(t)}
    =
    \sum_{j=1}^{n}w_{ij}s_j^{(t)}+\theta_i.
\]

\begin{figure}[ht]
    \centering
    \begin{tikzpicture}[
    neuron/.style={circle, draw=darkblue, very thick, fill=darkblue!10, minimum size=11mm},
    conn/.style={thick, draw=darkblue!70}
]
    \foreach \ang/\name/\lab in {90/v1/\(s_1\),18/v2/\(s_2\),-54/v3/\(s_3\),-126/v4/\(s_4\),162/v5/\(s_5\)}
        \node[neuron] (\name) at (\ang:2.1cm) {\lab};
    \foreach \a/\b in {v1/v2,v1/v3,v1/v4,v1/v5,v2/v3,v2/v4,v2/v5,v3/v4,v3/v5,v4/v5}
        \draw[conn] (\a) -- (\b);
    \node[fill=white, inner sep=2mm] at (0,0) {\(w_{ij}=w_{ji}\)};
\end{tikzpicture}

    \caption{Hopfieldova síť má symetrické spoje a~stav uložený v~neuronech.}
    \label{fig:ai-hopfield-sit}
\end{figure}

Při \emph{asynchronní aktualizaci} vybereme jediný neuron \(i\) a~nastavíme
\[
    s_i^{(t+1)}
    =
    \begin{cases}
        1, & h_i^{(t)}>0,\\
        -1, & h_i^{(t)}<0,\\
        s_i^{(t)}, & h_i^{(t)}=0.
    \end{cases}
\]
Ostatní složky stavu zůstanou nezměněné. Neurony můžeme procházet v~pevném nebo náhodném pořadí. Při synchronní aktualizaci bychom měnili všechny neurony najednou; ta však může mezi několika stavy kmitat.

Asynchronní pravidlo používá vždy nejnovější dostupný stav: po změně jednoho neuronu už následující neuron pracuje s~touto novou hodnotou. Pořadí aktualizací může ovlivnit, do kterého stabilního stavu síť dospěje, ale při symetrických vahách neporuší vlastnost poklesu energie dokázanou níže.

\subsection{Asociativní paměť}

Hopfieldova síť slouží jako \emph{asociativní paměť}. Neuchovává záznam pod samostatnou adresou. Zapamatovaný vzor je stabilní stav sítě a~je možné jej vybavit i~z~neúplné nebo mírně poškozené podoby. Opakované aktualizace postupně opravují jednotlivé složky, dokud síť nedospěje ke stabilnímu stavu.

\begin{figure}[ht]
    \centering
    \begin{tikzpicture}[
    state/.style={draw=darkblue, rounded corners, fill=darkblue!8, minimum width=29mm, minimum height=15mm, align=center},
    conn/.style={-Latex, very thick, draw=darkblue}
]
    \node[state] (damaged) at (0,0) {poškozený stav\\\(\mathbf{s}^{(0)}\)};
    \node[state] (middle) at (4.7,0) {opakované\\aktualizace};
    \node[state] (stored) at (9.4,0) {uložený vzor\\\(\mathbf{x}_r\)};
    \draw[conn] (damaged) -- node[above] {doplnění} (middle);
    \draw[conn] (middle) -- node[above] {ustálení} (stored);
    \node[below=7mm of middle] {nejbližší energetické minimum};
\end{tikzpicture}

    \caption{Vybavení vzoru z~jeho poškozené podoby.}
    \label{fig:ai-hopfield-vybaveni}
\end{figure}

Abychom odlišili počet neuronů od počtu ukládaných vzorů, zavedeme nové označení \(m\) pro počet vzorů. Nechť
\[
    \mathbf{x}_1,\ldots,\mathbf{x}_m
    \in\set{-1,1}^{n}
\]
jsou vzory určené k~uložení. Jednoduché Hebbovo pravidlo nastaví pro \(i\neq j\)
\[
    w_{ij}
    =
    \frac1n
    \sum_{r=1}^{m}x_{ri}x_{rj},
    \qquad
    w_{ii}=0,
\]
kde \(x_{ri}\) je \(i\)-tá složka vzoru \(\mathbf{x}_r\). Shodují-li se dvě složky ve většině vzorů, vznikne mezi nimi kladná váha; bývají-li opačné, váha je záporná.

Učení zde tedy neprobíhá opakovaným gradientním sestupem. Váhy vzniknou přímo jedním součtem přes ukládané vzory a~potom se při vybavování již nemění. Opakovaně se mění pouze stav neuronů, dokud síť nenalezne stabilní konfiguraci.

\subsection{Energetická funkce}

Stavu \(\mathbf{s}\in\set{-1,1}^{n}\) přiřadíme energii
\[
    E(\mathbf{s})
    =
    -\frac12
    \sum_{i=1}^{n}\sum_{j=1}^{n}
    w_{ij}s_is_j
    -
    \sum_{i=1}^{n}\theta_i s_i.
\]
Změní-li se asynchronně pouze neuron \(i\) ze stavu \(s_i\) do stavu \(s_i'\), symetrie vah dává
\[
    E(\mathbf{s}')
    -
    E(\mathbf{s})
    =
    -(s_i'-s_i)
    \left(
        \sum_{j=1}^{n}w_{ij}s_j+\theta_i
    \right)
    =
    -(s_i'-s_i)h_i.
\]
Když se neuron podle uvedeného pravidla skutečně změní, mají \(s_i'-s_i\) a~\(h_i\) stejné znaménko. Proto
\[
    E(\mathbf{s}')-E(\mathbf{s})<0.
\]
Při každé změně tedy energie přísně klesne. Síť má jen \(2^n\) stavů, a~proto nemůže měnit stav donekonečna. Po konečném počtu změn dospěje k~lokálnímu minimu energie, tedy ke stabilnímu stavu.

Tento důkaz se opírá o~tři podstatné předpoklady: váhy jsou symetrické, na diagonále platí \(w_{ii}=0\) a~v~jednom kroku se mění pouze jeden neuron. U~synchronní aktualizace se mohou změny různých neuronů navzájem ovlivnit a~síť může přecházet mezi dvěma stavy, takže stejný argument nelze bez úpravy použít.

\begin{center}
    \centering
    \begin{tikzpicture}
    \begin{axis}[
        width=0.78\textwidth,
        height=5.3cm,
        axis lines=left,
        xlabel={aktualizace},
        ylabel={\(E(\mathbf{s})\)},
        xmin=0, xmax=8.6,
        ymin=0, ymax=6.7,
        xtick=\empty,
        ytick=\empty,
        clip=false
    ]
        \addplot[darkblue, very thick, const plot] coordinates {
            (0,6.0) (1,5.1) (2,4.2) (3,4.2) (4,2.8)
            (5,2.0) (6,2.0) (7,1.2) (8,1.2)
        };
        \node[above] at (axis cs:7.5,1.55) {stabilní stav};
        \draw[-Latex, thick] (axis cs:7.5,1.45) -- (axis cs:7.5,1.23);
    \end{axis}
\end{tikzpicture}

    \captionof{figure}{Aktualizace snižují energii, dokud síť nedospěje ke stabilnímu stavu.}
    \label{fig:ai-hopfield-energie}
\end{center}

\subsection{Implementace a~omezení}

Program~\ref{lst:ai-hopfield} vytvoří váhy Hebbovým pravidlem a~poškozený vstup aktualizuje asynchronně. Nulové lokální pole ponechá původní stav neuronu, což odpovídá pravidlu použitému v~důkazu poklesu energie.

Pseudokód odděluje dvě fáze. Procedura učení nejprve ze vzorů vytvoří symetrickou matici vah. Procedura vybavení potom přijme počáteční stav a~provádí celé průchody neurony tak dlouho, dokud se v~jednom průchodu žádný stav nezmění nebo dokud nedosáhne bezpečnostního limitu.

\begin{algorithm}
    \caption{\textsc{hopfield-recall}}
    \label{alg:ai-hopfield}
    \KwIn{Vzory \(\mathbf{x}_1,\ldots,\mathbf{x}_m\in\set{-1,1}^n\), prahy \(\theta_i\), počáteční stav \(\mathbf{s}\) a~nejvyšší počet průchodů \(T\)}
    \For{\(i\gets1,2,\ldots,n\)}{
        \For{\(j\gets1,2,\ldots,n\)}{
            \If{\(i=j\)}{\(w_{ij}\gets0\)}
            \If{\(i\neq j\)}{\(\displaystyle w_{ij}\gets\frac1n\sum_{q=1}^{m}x_{qi}x_{qj}\)}
        }
    }
    \For{\(r\gets1,2,\ldots,T\)}{
        \(zmena\gets0\)\\
        \For{\(i\gets1,2,\ldots,n\)}{
            \(\displaystyle h_i\gets\sum_{j=1}^{n}w_{ij}s_j+\theta_i\)\\
            \If{\(h_i>0\)}{\(s_i'\gets1\)}
            \If{\(h_i<0\)}{\(s_i'\gets-1\)}
            \If{\(h_i=0\)}{\(s_i'\gets s_i\)}
            \If{\(s_i'\neq s_i\)}{
                \(s_i\gets s_i'\), \(zmena\gets1\)
            }
        }
        \lIf{\(zmena=0\)}{\Return{\(\mathbf{s}\)}}
    }
    \KwOut{Poslední stav \(\mathbf{s}\)}
\end{algorithm}

\lstinputlisting[
    caption={Hopfieldova asociativní paměť.},
    label={lst:ai-hopfield}
]{07-ai/code/hopfield.py}

Kapacita Hopfieldovy sítě je omezená. Při příliš velkém počtu uložených nebo silně podobných vzorů se jejich příspěvky ve vahách ruší. Síť potom může skončit ve špatném uloženém vzoru nebo v~nežádoucím stabilním stavu, který nebyl součástí trénovacích dat. Model je proto užitečný hlavně jako názorný příklad rekurentní sítě, asociativní paměti a~učení bez učitele.

Kód záměrně používá malé binární vzory, aby bylo možné sledovat každý krok. U~obrázků bychom pixely nejprve převedli na hodnoty \(-1\) a~\(1\), obraz zploštili do vektoru a~po vybavení jej opět složili do původního tvaru. Princip vah i~aktualizací by zůstal beze změny.

\begin{remark}
    John Hopfield publikoval známý model asociativní paměti roku 1982. Jeho práce propojila dřívější myšlenky rekurentních sítí s~energetickou funkcí známou ze statistické fyziky a~výrazně přispěla k~obnovenému zájmu o~neuronové sítě.
\end{remark}
